%======================================================================
\section{Related Work}
\label{sec:related}
%======================================================================

\subsection{LLM-Assisted Vulnerability Detection}

Recent work has explored using LLMs to \emph{detect} vulnerabilities in smart contracts, complementing traditional tools. GPTScan~\cite{gptscan} combines GPT-4 with static analysis to identify logic vulnerabilities, achieving higher precision than standalone tools on a benchmark of 168 contracts. SmartLLMSentry~\cite{smartllmsentry} employs LLMs as vulnerability detectors with chain-of-thought reasoning. PropertyGPT~\cite{propertygpt} uses LLMs to generate formal security properties for automated verification. David et al.~\cite{david2023you} evaluated ChatGPT's ability to detect and repair 52 smart contract vulnerabilities, finding moderate detection accuracy (68\%) but low repair success rates. ContractTinker~\cite{contracttinker} proposes LLM-based automated smart contract repair.

These works are complementary to ours: they study LLMs as \emph{detectors}, while we study LLMs as vulnerability \emph{sources}. An important implication of combining both perspectives is the potential for circular vulnerabilities---LLMs that generate vulnerable code may also fail to detect the same vulnerability classes in their output.

\subsection{Smart Contract Security Analysis}

The smart contract security tooling ecosystem has matured significantly. Slither~\cite{slither} provides fast static analysis with over 90 built-in detectors, making it the most widely used tool for Solidity auditing. Mythril~\cite{mythril} performs symbolic execution for deep bug detection, including those requiring specific state conditions. Securify~\cite{securify} applies formal verification to check compliance with predefined security patterns. Oyente~\cite{luu2016making} was among the first symbolic execution tools for Ethereum smart contracts. Manticore~\cite{mossberg2019manticore} provides advanced symbolic analysis with support for multiple architectures. SmartCheck~\cite{smartcheck} uses structural pattern-based detection.

Durieux et al.~\cite{durieux2020empirical} conducted the most comprehensive empirical evaluation of smart contract analysis tools, testing 9 tools on 47,587 contracts. Their key finding---that no single tool achieves complete coverage and that tools have low inter-agreement---directly motivated our multi-tool approach. Ren et al.~\cite{ren2021empirical} conducted an empirical study of vulnerabilities in 21,212 real-world contracts, establishing prevalence baselines. SmartBugs~\cite{smartbugs} provides a curated, labeled benchmark dataset for tool evaluation. We build on these tools as our analysis infrastructure and contribute a new dimension by applying them to LLM-generated rather than human-written code.

\subsection{Security of LLM-Generated Code}

The security implications of LLM code generation have been studied primarily in general-purpose languages. Pearce et al.~\cite{copilot_eval} conducted the foundational study evaluating GitHub Copilot's security across 89 scenarios derived from CWE types, finding that approximately 40\% of generated suggestions contained vulnerabilities. Perry et al.~\cite{perry2023users} demonstrated through a controlled study with 47 participants that developers using AI assistants produce less secure code while believing it to be more secure. Sandoval et al.~\cite{sandoval2023lost} confirmed these findings in a larger study with professional developers.

Tony et al.~\cite{tony2023llmSecEval} created LLMSecEval, a benchmark of 150 test cases for evaluating code LLM security across 18 CWE types. Tihanyi et al.~\cite{tihanyi2023formai} proposed FormAI, a dataset of 112,000 AI-generated C programs with vulnerability labels obtained through dynamic analysis. Jesse et al.~\cite{jesse2023large} studied LLM-generated code bugs more broadly, finding systematic patterns in error generation.

Our work extends this line of research to the smart contract domain, which introduces unique challenges: immutability prevents post-deployment patching, financial exposure creates immediate exploitation incentives, and the adversarial blockchain environment means any vulnerability is publicly discoverable. Furthermore, our adversarial prompt framework and cross-LLM comparison methodology are novel contributions not present in prior work.

\subsection{Adversarial Attacks on LLMs}

Prompt injection~\cite{perez2022ignore, greshake2023not} and jailbreaking~\cite{kang2024exploiting} attacks demonstrate that LLMs can be manipulated to produce harmful or unintended outputs. In the code domain, Schuster et al.~\cite{schuster2021you} showed that training data poisoning can introduce targeted vulnerabilities into code completion models. Wan et al.~\cite{wan2023poisoning} extended this to instruction-tuned code generation models, demonstrating that subtle training data manipulation can cause models to systematically produce specific vulnerability types.

Our adversarial prompt framework differs from these prior approaches in two ways: (1)~we focus on \emph{prompt-level} manipulation rather than training data poisoning, reflecting a more realistic threat model; and (2)~our adversarial prompts are designed to appear natural and plausible---resembling requests a non-security-aware developer might genuinely write---rather than explicit jailbreaking attempts. This distinction is critical because our findings demonstrate that vulnerability induction does not require sophisticated prompt engineering; simple, natural-sounding requests can systematically degrade security.

\subsection{Smart Contract Datasets and Benchmarks}

Several datasets exist for smart contract security research. SmartBugs~\cite{smartbugs} provides a curated dataset of vulnerable contracts for tool benchmarking, categorized by vulnerability type. The SWC Registry~\cite{swc_registry} maintains a taxonomy of known smart contract weaknesses with test cases. Ren et al.~\cite{ren2021empirical} compiled a dataset of 21,212 real-world contracts with vulnerability annotations from multiple tools.

Our dataset differs fundamentally from these prior datasets: it captures vulnerabilities \emph{introduced by LLM generation} rather than curated from known bugs, mined from deployed contracts, or synthetically constructed. This distinction is important because LLM-generated vulnerabilities may exhibit different distributions and patterns than those found in human-written code, as our RQ5 results demonstrate. Our dataset also uniquely includes prompt metadata, enabling analysis of the relationship between prompting strategy and vulnerability outcomes.
